% Auto-generated full source appendix. Regenerate: python scripts/gen_source_appendix.py

\subsection{\texttt{src/config.py}}
Configuration and hyperparameters.
\begin{lstlisting}[caption={Configuration and hyperparameters (\texttt{src/config.py})}]
"""Central configuration for the acoustic scene awareness project.

Every tunable lives here so experiments stay reproducible and the report can
quote exact settings.
"""
from pathlib import Path

import torch

# --------------------------------------------------------------------------
# Paths
# --------------------------------------------------------------------------
ROOT = Path(__file__).resolve().parent.parent

RAW_DIR = ROOT / "data" / "raw"
PROCESSED_DIR = ROOT / "data" / "processed"
MODELS_DIR = ROOT / "models"
RESULTS_DIR = ROOT / "results"
FIGURES_DIR = RESULTS_DIR / "figures"
LOGS_DIR = RESULTS_DIR / "logs"

US8K_DIR = RAW_DIR / "UrbanSound8K"
US8K_AUDIO = US8K_DIR / "audio"
US8K_META = US8K_DIR / "metadata" / "UrbanSound8K.csv"

ESC50_DIR = RAW_DIR / "ESC-50-master"
ESC50_AUDIO = ESC50_DIR / "audio"
ESC50_META = ESC50_DIR / "meta" / "esc50.csv"

for _d in (PROCESSED_DIR, MODELS_DIR, FIGURES_DIR, LOGS_DIR):
    _d.mkdir(parents=True, exist_ok=True)

# --------------------------------------------------------------------------
# Audio / mel-spectrogram front-end
# --------------------------------------------------------------------------
# UrbanSound8K clips are up to 4 s. Fixing every clip to 4 s at 22.05 kHz keeps
# the spectrogram a constant shape, which the CNN requires.
SAMPLE_RATE = 22050
CLIP_SECONDS = 4.0
N_SAMPLES = int(SAMPLE_RATE * CLIP_SECONDS)

N_FFT = 1024
HOP_LENGTH = 512
N_MELS = 64
F_MIN = 20
F_MAX = SAMPLE_RATE // 2

# Resulting spectrogram shape -> (N_MELS, N_FRAMES)
N_FRAMES = int(N_SAMPLES / HOP_LENGTH) + 1

# --------------------------------------------------------------------------
# Training
# --------------------------------------------------------------------------
BATCH_SIZE = 64
EPOCHS = 40
LEARNING_RATE = 1e-3
WEIGHT_DECAY = 1e-4
EARLY_STOP_PATIENCE = 8
SEED = 42

# UrbanSound8K ships 10 predefined folds. The official protocol is 10-fold
# cross-validation using these folds — never a random split, because clips cut
# from the same source recording live in the same fold. Random splitting leaks
# them across train/test and inflates accuracy.
N_FOLDS = 10

# For fast iteration we can evaluate on a subset of folds; the full run uses all.
QUICK_FOLDS = [1]

# --------------------------------------------------------------------------
# Device — Apple Silicon GPU via Metal (MPS), else CPU
# --------------------------------------------------------------------------
def get_device() -> torch.device:
    if torch.backends.mps.is_available():
        return torch.device("mps")
    if torch.cuda.is_available():
        return torch.device("cuda")
    return torch.device("cpu")


DEVICE = get_device()

# --------------------------------------------------------------------------
# Classes
# --------------------------------------------------------------------------
US8K_CLASSES = [
    "air_conditioner",
    "car_horn",
    "children_playing",
    "dog_bark",
    "drilling",
    "engine_idling",
    "gun_shot",
    "jackhammer",
    "siren",
    "street_music",
]

# Human-facing labels for the live demo, plus which sounds are treated as
# safety-critical alerts for the assistive use case.
DEMO_LABELS = {
    "air_conditioner": "Air conditioner",
    "car_horn": "Car horn",
    "children_playing": "Children playing",
    "dog_bark": "Dog barking",
    "drilling": "Drilling",
    "engine_idling": "Engine idling",
    "gun_shot": "Gunshot",
    "jackhammer": "Jackhammer",
    "siren": "Siren",
    "street_music": "Street music",
}

ALERT_CLASSES = {"car_horn", "gun_shot", "siren", "dog_bark"}

# ESC-50 classes used for the novel-sound (anomaly) experiment. These are
# household/safety sounds the UrbanSound8K model has never been trained on, so
# a good novelty detector should flag them as unknown.
ESC50_NOVEL_CLASSES = [
    "door_wood_knock",
    "glass_breaking",
    "crying_baby",
    "clock_alarm",
    "door_wood_creaks",
    "water_drops",
    "footsteps",
    "can_opening",
]
\end{lstlisting}

\subsection{\texttt{src/data.py}}
Data pipeline: loading, mel-spectrograms, augmentation.
\begin{lstlisting}[caption={Data pipeline: loading, mel-spectrograms, augmentation (\texttt{src/data.py})}]
"""Audio loading, mel-spectrogram feature extraction and PyTorch datasets.

The expensive part of this project is turning ~10k wav files into
mel-spectrograms. That is done once by `precompute_features()` and cached to
.npy, after which training reads features straight from memory.
"""
from __future__ import annotations

import numpy as np
import pandas as pd
import torch
from torch.utils.data import Dataset

from . import config as C


# ---------------------------------------------------------------------------
# Metadata
# ---------------------------------------------------------------------------
def load_us8k_metadata() -> pd.DataFrame:
    """UrbanSound8K metadata, one row per clip, including its official fold."""
    df = pd.read_csv(C.US8K_META)
    df["path"] = df.apply(
        lambda r: C.US8K_AUDIO / f"fold{r['fold']}" / r["slice_file_name"], axis=1
    )
    return df


def load_esc50_metadata() -> pd.DataFrame:
    """ESC-50 metadata. Used only as a source of *unseen* sounds for the
    novelty-detection experiment."""
    df = pd.read_csv(C.ESC50_META)
    df["path"] = df["filename"].apply(lambda f: C.ESC50_AUDIO / f)
    return df


# ---------------------------------------------------------------------------
# Feature extraction
# ---------------------------------------------------------------------------
def load_audio(path, sr: int = C.SAMPLE_RATE) -> np.ndarray:
    """Load a clip as mono at `sr`, then pad/trim to exactly CLIP_SECONDS.

    Fixing the length is what makes every spectrogram the same shape, which the
    CNN needs. Clips shorter than 4 s are zero-padded; longer ones are cut.
    """
    import librosa

    y, _ = librosa.load(path, sr=sr, mono=True)
    if len(y) < C.N_SAMPLES:
        y = np.pad(y, (0, C.N_SAMPLES - len(y)), mode="constant")
    else:
        y = y[: C.N_SAMPLES]
    return y


def audio_to_melspec(y: np.ndarray, sr: int = C.SAMPLE_RATE) -> np.ndarray:
    """Convert a waveform to a log-scaled mel-spectrogram.

    This is the core "audio becomes an image" step: the result is a
    (N_MELS, N_FRAMES) matrix of energy per frequency band over time, which a
    2-D CNN can process exactly like a grayscale image.

    Decibel scaling matters — human loudness perception is logarithmic, and
    without it a handful of loud frames dominate the input range.
    """
    import librosa

    mel = librosa.feature.melspectrogram(
        y=y,
        sr=sr,
        n_fft=C.N_FFT,
        hop_length=C.HOP_LENGTH,
        n_mels=C.N_MELS,
        fmin=C.F_MIN,
        fmax=C.F_MAX,
    )
    mel_db = librosa.power_to_db(mel, ref=np.max)
    return mel_db.astype(np.float32)


def extract_features(path) -> np.ndarray:
    return audio_to_melspec(load_audio(path))


def precompute_features(df: pd.DataFrame, cache_name: str, force: bool = False):
    """Extract mel-spectrograms for every row of `df` and cache them.

    Returns (X, meta) where X has shape (N, N_MELS, N_FRAMES).
    """
    from tqdm import tqdm

    cache_x = C.PROCESSED_DIR / f"{cache_name}_X.npy"
    cache_meta = C.PROCESSED_DIR / f"{cache_name}_meta.csv"

    if cache_x.exists() and cache_meta.exists() and not force:
        X = np.load(cache_x)
        meta = pd.read_csv(cache_meta)
        print(f"[cache] loaded {cache_name}: X={X.shape}")
        return X, meta

    feats, keep = [], []
    for i, row in tqdm(
        df.iterrows(), total=len(df), desc=f"extracting {cache_name}", unit="clip"
    ):
        try:
            feats.append(extract_features(row["path"]))
            keep.append(i)
        except Exception as exc:  # a handful of US8K files are known to be odd
            print(f"  ! skipped {row['path']}: {exc}")

    X = np.stack(feats)
    meta = df.loc[keep].drop(columns=["path"]).reset_index(drop=True)

    np.save(cache_x, X)
    meta.to_csv(cache_meta, index=False)
    print(f"[cache] saved {cache_name}: X={X.shape} -> {cache_x}")
    return X, meta


# ---------------------------------------------------------------------------
# Normalisation
# ---------------------------------------------------------------------------
def compute_norm_stats(X: np.ndarray) -> tuple[float, float]:
    """Mean/std computed on the TRAIN split only, then applied to val/test.

    Computing these over the whole dataset would leak test statistics into
    training, so they are always derived from the training fold.
    """
    return float(X.mean()), float(X.std() + 1e-8)


# ---------------------------------------------------------------------------
# Augmentation (spectrogram domain)
# ---------------------------------------------------------------------------
def spec_augment(
    spec: np.ndarray,
    rng: np.random.Generator,
    n_freq_masks: int = 1,
    n_time_masks: int = 1,
    freq_mask_max: int = 8,
    time_mask_max: int = 16,
) -> np.ndarray:
    """SpecAugment: blank out random frequency bands and time spans.

    Forces the network to use the whole time-frequency pattern instead of
    latching onto one narrow band, which measurably reduces overfitting on a
    dataset this small.
    """
    spec = spec.copy()
    n_mels, n_frames = spec.shape
    fill = spec.min()

    for _ in range(n_freq_masks):
        f = int(rng.integers(0, freq_mask_max + 1))
        if f > 0 and n_mels > f:
            f0 = int(rng.integers(0, n_mels - f))
            spec[f0 : f0 + f, :] = fill

    for _ in range(n_time_masks):
        t = int(rng.integers(0, time_mask_max + 1))
        if t > 0 and n_frames > t:
            t0 = int(rng.integers(0, n_frames - t))
            spec[:, t0 : t0 + t] = fill

    return spec


def time_shift(spec: np.ndarray, rng: np.random.Generator, max_frac: float = 0.2):
    """Roll the clip in time — a siren is a siren regardless of when it starts."""
    n_frames = spec.shape[1]
    shift = int(rng.integers(-int(n_frames * max_frac), int(n_frames * max_frac) + 1))
    return np.roll(spec, shift, axis=1)


# ---------------------------------------------------------------------------
# Datasets
# ---------------------------------------------------------------------------
class SpectrogramDataset(Dataset):
    """Serves normalised (1, N_MELS, N_FRAMES) tensors and integer labels."""

    def __init__(
        self,
        X: np.ndarray,
        y: np.ndarray,
        mean: float,
        std: float,
        augment: bool = False,
        seed: int = C.SEED,
    ):
        self.X = X
        self.y = y.astype(np.int64)
        self.mean = mean
        self.std = std
        self.augment = augment
        self.rng = np.random.default_rng(seed)

    def __len__(self) -> int:
        return len(self.X)

    def __getitem__(self, idx: int):
        spec = self.X[idx]
        if self.augment:
            spec = time_shift(spec, self.rng)
            spec = spec_augment(spec, self.rng)
        spec = (spec - self.mean) / self.std
        return torch.from_numpy(spec).unsqueeze(0).float(), int(self.y[idx])


def make_fold_split(X, meta, test_fold: int, val_fold: int | None = None):
    """Split by UrbanSound8K's official folds.

    Clips sliced from the same source recording share a fold, so splitting by
    fold (rather than randomly) is what keeps the evaluation honest.
    """
    if val_fold is None:
        val_fold = test_fold % C.N_FOLDS + 1

    folds = meta["fold"].values
    labels = meta["classID"].values

    te = folds == test_fold
    va = folds == val_fold
    tr = ~(te | va)

    return (X[tr], labels[tr]), (X[va], labels[va]), (X[te], labels[te])
\end{lstlisting}

\subsection{\texttt{src/models.py}}
Network architectures: MLP, CNN, autoencoder.
\begin{lstlisting}[caption={Network architectures: MLP, CNN, autoencoder (\texttt{src/models.py})}]
"""Neural network architectures.

Three networks, each serving a distinct task in the project:

  MLP           baseline classifier — flattens the spectrogram, throwing away
                all time-frequency structure. Exists to quantify what the CNN's
                inductive bias is actually worth.
  AudioCNN      main classifier — 2-D convolutions over the mel-spectrogram.
  ConvAutoencoder  novelty detector — trained only to reconstruct known sounds,
                so unfamiliar sounds reconstruct badly and can be flagged.
"""
from __future__ import annotations

import torch
import torch.nn as nn
import torch.nn.functional as F

from . import config as C


# ---------------------------------------------------------------------------
# Baseline: multi-layer perceptron
# ---------------------------------------------------------------------------
class MLPBaseline(nn.Module):
    """A fully-connected net over the flattened spectrogram.

    Every input pixel gets its own weight, so the model has no notion that two
    neighbouring time frames are related. It also carries far more parameters
    than the CNN despite being the weaker model — the comparison the report
    uses to motivate convolution.
    """

    def __init__(
        self,
        n_mels: int = C.N_MELS,
        n_frames: int = C.N_FRAMES,
        n_classes: int = 10,
        hidden: tuple[int, ...] = (512, 256),
        dropout: float = 0.5,
    ):
        super().__init__()
        in_dim = n_mels * n_frames

        layers: list[nn.Module] = [nn.Flatten()]
        prev = in_dim
        for h in hidden:
            layers += [nn.Linear(prev, h), nn.BatchNorm1d(h), nn.ReLU(), nn.Dropout(dropout)]
            prev = h
        layers.append(nn.Linear(prev, n_classes))

        self.net = nn.Sequential(*layers)

    def forward(self, x):
        return self.net(x)


# ---------------------------------------------------------------------------
# Main classifier: convolutional neural network
# ---------------------------------------------------------------------------
class ConvBlock(nn.Module):
    """Conv -> BatchNorm -> ReLU -> MaxPool, the standard CNN building block."""

    def __init__(self, in_ch: int, out_ch: int, pool: int = 2, dropout: float = 0.0):
        super().__init__()
        self.conv = nn.Conv2d(in_ch, out_ch, kernel_size=3, padding=1, bias=False)
        self.bn = nn.BatchNorm2d(out_ch)
        self.pool = nn.MaxPool2d(pool)
        self.drop = nn.Dropout2d(dropout) if dropout > 0 else nn.Identity()

    def forward(self, x):
        x = F.relu(self.bn(self.conv(x)))
        x = self.pool(x)
        return self.drop(x)


class AudioCNN(nn.Module):
    """Four-block CNN over log-mel spectrograms.

    Early filters learn local time-frequency texture (onsets, harmonic stacks,
    broadband noise); deeper blocks compose those into class-level patterns.

    Global average pooling replaces a large flatten+dense head. That keeps the
    parameter count ~10x below the MLP and makes the model robust to a sound
    occurring at a different point in the clip.

    `return_embedding=True` exposes the 256-d penultimate features, which the
    clustering / t-SNE analysis uses to show what the network learned.
    """

    def __init__(self, n_classes: int = 10, dropout: float = 0.3, width: int = 32):
        super().__init__()
        w = width
        self.block1 = ConvBlock(1, w, pool=2)
        self.block2 = ConvBlock(w, w * 2, pool=2, dropout=dropout / 2)
        self.block3 = ConvBlock(w * 2, w * 4, pool=2, dropout=dropout / 2)
        self.block4 = ConvBlock(w * 4, w * 8, pool=2, dropout=dropout)

        self.gap = nn.AdaptiveAvgPool2d(1)
        self.dropout = nn.Dropout(dropout)
        self.fc = nn.Linear(w * 8, n_classes)
        self.embed_dim = w * 8

    def forward(self, x, return_embedding: bool = False):
        x = self.block1(x)
        x = self.block2(x)
        x = self.block3(x)
        x = self.block4(x)

        emb = self.gap(x).flatten(1)
        logits = self.fc(self.dropout(emb))

        if return_embedding:
            return logits, emb
        return logits


# ---------------------------------------------------------------------------
# Novelty detection: convolutional autoencoder
# ---------------------------------------------------------------------------
class ConvAutoencoder(nn.Module):
    """Compresses a spectrogram to a small latent code and rebuilds it.

    Trained only on the ten known UrbanSound8K classes, it learns to reconstruct
    those well. A sound it has never encountered (a smashing window, a crying
    baby) falls outside that learned manifold and reconstructs poorly, so the
    per-clip reconstruction error works directly as a novelty score.
    """

    def __init__(self, latent_ch: int = 32):
        super().__init__()

        self.encoder = nn.Sequential(
            nn.Conv2d(1, 16, 3, stride=2, padding=1),   # -> 32 x 87
            nn.BatchNorm2d(16),
            nn.ReLU(),
            nn.Conv2d(16, 32, 3, stride=2, padding=1),  # -> 16 x 44
            nn.BatchNorm2d(32),
            nn.ReLU(),
            nn.Conv2d(32, latent_ch, 3, stride=2, padding=1),  # -> 8 x 22
            nn.BatchNorm2d(latent_ch),
            nn.ReLU(),
        )

        self.decoder = nn.Sequential(
            nn.ConvTranspose2d(latent_ch, 32, 3, stride=2, padding=1, output_padding=1),
            nn.BatchNorm2d(32),
            nn.ReLU(),
            nn.ConvTranspose2d(32, 16, 3, stride=2, padding=1, output_padding=1),
            nn.BatchNorm2d(16),
            nn.ReLU(),
            nn.ConvTranspose2d(16, 1, 3, stride=2, padding=1, output_padding=1),
        )

    def forward(self, x):
        z = self.encoder(x)
        out = self.decoder(z)
        # Transposed convolutions can overshoot by a pixel; crop to match input.
        return out[..., : x.shape[-2], : x.shape[-1]]

    def encode(self, x):
        return self.encoder(x)

    @torch.no_grad()
    def reconstruction_error(self, x) -> torch.Tensor:
        """Per-sample mean squared error — the novelty score."""
        recon = self(x)
        return ((recon - x) ** 2).flatten(1).mean(dim=1)


# ---------------------------------------------------------------------------
# Helpers
# ---------------------------------------------------------------------------
def count_parameters(model: nn.Module) -> int:
    return sum(p.numel() for p in model.parameters() if p.requires_grad)


def build_model(name: str, n_classes: int = 10, **kwargs) -> nn.Module:
    name = name.lower()
    if name == "mlp":
        return MLPBaseline(n_classes=n_classes, **kwargs)
    if name == "cnn":
        return AudioCNN(n_classes=n_classes, **kwargs)
    if name in ("autoencoder", "ae"):
        return ConvAutoencoder(**kwargs)
    raise ValueError(f"unknown model: {name}")
\end{lstlisting}

\subsection{\texttt{src/train.py}}
Training loop and 10-fold cross-validation.
\begin{lstlisting}[caption={Training loop and 10-fold cross-validation (\texttt{src/train.py})}]
"""Training engine: single-fold training plus the official 10-fold protocol."""
from __future__ import annotations

import json
import time
from dataclasses import dataclass, field, asdict

import numpy as np
import torch
import torch.nn as nn
from torch.utils.data import DataLoader

from . import config as C
from .data import SpectrogramDataset, compute_norm_stats, make_fold_split
from .models import build_model, count_parameters


def set_seed(seed: int = C.SEED) -> None:
    import random

    random.seed(seed)
    np.random.seed(seed)
    torch.manual_seed(seed)
    torch.mps.manual_seed(seed) if torch.backends.mps.is_available() else None


@dataclass
class FoldResult:
    fold: int
    model: str
    best_val_acc: float
    test_acc: float
    test_f1_macro: float
    epochs_run: int
    train_seconds: float
    n_params: int
    history: dict = field(default_factory=dict)


# ---------------------------------------------------------------------------
# Epoch loops
# ---------------------------------------------------------------------------
def _run_epoch(model, loader, criterion, optimizer, device, train: bool):
    model.train() if train else model.eval()
    total_loss, correct, total = 0.0, 0, 0

    with torch.set_grad_enabled(train):
        for xb, yb in loader:
            xb, yb = xb.to(device), yb.to(device)

            logits = model(xb)
            loss = criterion(logits, yb)

            if train:
                optimizer.zero_grad(set_to_none=True)
                loss.backward()
                optimizer.step()

            total_loss += loss.item() * xb.size(0)
            correct += (logits.argmax(1) == yb).sum().item()
            total += xb.size(0)

    return total_loss / total, correct / total


@torch.no_grad()
def predict(model, loader, device):
    """Returns (y_true, y_pred, probabilities) for a whole loader."""
    model.eval()
    ys, ps, probs = [], [], []
    for xb, yb in loader:
        logits = model(xb.to(device))
        p = torch.softmax(logits, dim=1)
        ys.append(yb.numpy())
        ps.append(logits.argmax(1).cpu().numpy())
        probs.append(p.cpu().numpy())
    return np.concatenate(ys), np.concatenate(ps), np.concatenate(probs)


# ---------------------------------------------------------------------------
# Single fold
# ---------------------------------------------------------------------------
def train_one_fold(
    X,
    meta,
    test_fold: int,
    model_name: str = "cnn",
    epochs: int = C.EPOCHS,
    batch_size: int = C.BATCH_SIZE,
    lr: float = C.LEARNING_RATE,
    augment: bool = True,
    device: torch.device = C.DEVICE,
    n_classes: int = 10,
    verbose: bool = True,
    save_path=None,
) -> tuple[nn.Module, FoldResult]:
    """Train one model on the folds outside `test_fold`, evaluate on it."""
    from sklearn.metrics import f1_score

    set_seed(C.SEED + test_fold)

    (Xtr, ytr), (Xva, yva), (Xte, yte) = make_fold_split(X, meta, test_fold)

    # Normalisation statistics come from the training split only.
    mean, std = compute_norm_stats(Xtr)

    ds_tr = SpectrogramDataset(Xtr, ytr, mean, std, augment=augment)
    ds_va = SpectrogramDataset(Xva, yva, mean, std, augment=False)
    ds_te = SpectrogramDataset(Xte, yte, mean, std, augment=False)

    dl_tr = DataLoader(ds_tr, batch_size=batch_size, shuffle=True, drop_last=True)
    dl_va = DataLoader(ds_va, batch_size=batch_size)
    dl_te = DataLoader(ds_te, batch_size=batch_size)

    model = build_model(model_name, n_classes=n_classes).to(device)
    criterion = nn.CrossEntropyLoss(label_smoothing=0.05)
    optimizer = torch.optim.AdamW(model.parameters(), lr=lr, weight_decay=C.WEIGHT_DECAY)
    scheduler = torch.optim.lr_scheduler.CosineAnnealingLR(optimizer, T_max=epochs)

    history = {"train_loss": [], "train_acc": [], "val_loss": [], "val_acc": []}
    best_val, best_state, patience, t0 = 0.0, None, 0, time.time()
    epochs_run = 0

    for ep in range(1, epochs + 1):
        tr_loss, tr_acc = _run_epoch(model, dl_tr, criterion, optimizer, device, True)
        va_loss, va_acc = _run_epoch(model, dl_va, criterion, optimizer, device, False)
        scheduler.step()
        epochs_run = ep

        history["train_loss"].append(tr_loss)
        history["train_acc"].append(tr_acc)
        history["val_loss"].append(va_loss)
        history["val_acc"].append(va_acc)

        if va_acc > best_val:
            best_val = va_acc
            best_state = {k: v.detach().cpu().clone() for k, v in model.state_dict().items()}
            patience = 0
        else:
            patience += 1

        if verbose and (ep % 5 == 0 or ep == 1):
            print(
                f"  fold{test_fold} {model_name} ep{ep:3d}/{epochs} "
                f"train {tr_loss:.3f}/{tr_acc:.3f}  val {va_loss:.3f}/{va_acc:.3f}"
            )

        if patience >= C.EARLY_STOP_PATIENCE:
            if verbose:
                print(f"  fold{test_fold} early stop at epoch {ep}")
            break

    if best_state is not None:
        model.load_state_dict(best_state)

    y_true, y_pred, _ = predict(model, dl_te, device)
    test_acc = float((y_true == y_pred).mean())
    test_f1 = float(f1_score(y_true, y_pred, average="macro"))

    result = FoldResult(
        fold=test_fold,
        model=model_name,
        best_val_acc=best_val,
        test_acc=test_acc,
        test_f1_macro=test_f1,
        epochs_run=epochs_run,
        train_seconds=time.time() - t0,
        n_params=count_parameters(model),
        history=history,
    )

    if save_path:
        torch.save(
            {"state_dict": model.state_dict(), "mean": mean, "std": std,
             "model_name": model_name, "n_classes": n_classes},
            save_path,
        )

    if verbose:
        print(
            f"  fold{test_fold} {model_name} -> test acc {test_acc:.4f} "
            f"f1 {test_f1:.4f} ({result.train_seconds:.0f}s)"
        )

    return model, result


# ---------------------------------------------------------------------------
# Full cross-validation
# ---------------------------------------------------------------------------
def cross_validate(
    X,
    meta,
    model_name: str = "cnn",
    folds=None,
    epochs: int = C.EPOCHS,
    augment: bool = True,
    tag: str | None = None,
    **kwargs,
):
    """Run the official 10-fold protocol and report mean +/- std.

    Two conveniences for long unattended runs, both controlled by environment
    variables so they never affect the results:

    * ``COOLDOWN_SEC`` pauses between folds to let the GPU shed heat. Folds are
      independent, so this changes timing only, never the numbers.
    * ``RESUME=1`` skips a configuration whose cached JSON already covers the
      requested folds, so a restarted run does not redo finished work. Results
      are deterministic, so a skipped config is identical to a rerun.
    """
    import os

    folds = folds or list(range(1, C.N_FOLDS + 1))
    cooldown = float(os.environ.get("COOLDOWN_SEC", "0"))

    name = tag or f"{model_name}_aug{int(augment)}"
    out = C.LOGS_DIR / f"cv_{name}.json"

    if os.environ.get("RESUME") == "1" and out.exists():
        cached = json.loads(out.read_text())
        if set(cached.get("folds", [])) >= set(folds):
            print(f"[resume] {name} already complete "
                  f"({cached['acc_mean']*100:.2f}%) -- skipping")
            return cached

    results: list[FoldResult] = []

    print(f"\n{'='*70}\n{model_name.upper()}  |  folds={folds}  augment={augment}"
          f"{f'  cooldown={cooldown:.0f}s' if cooldown else ''}\n{'='*70}")

    for i, f in enumerate(folds):
        _, res = train_one_fold(
            X, meta, test_fold=f, model_name=model_name,
            epochs=epochs, augment=augment, **kwargs
        )
        results.append(res)

        if cooldown and i < len(folds) - 1:
            print(f"  [cooldown] pausing {cooldown:.0f}s to let the GPU cool "
                  f"(fold {i+1}/{len(folds)} done)")
            time.sleep(cooldown)

    accs = np.array([r.test_acc for r in results])
    f1s = np.array([r.test_f1_macro for r in results])

    summary = {
        "model": model_name,
        "augment": augment,
        "folds": folds,
        "acc_mean": float(accs.mean()), "acc_std": float(accs.std()),
        "f1_mean": float(f1s.mean()), "f1_std": float(f1s.std()),
        "n_params": results[0].n_params,
        "total_seconds": float(sum(r.train_seconds for r in results)),
        "per_fold": [asdict(r) for r in results],
    }

    print(
        f"\n{model_name.upper()} SUMMARY  acc {accs.mean():.4f} +/- {accs.std():.4f} | "
        f"macro-F1 {f1s.mean():.4f} +/- {f1s.std():.4f} | params {results[0].n_params:,}"
    )

    name = tag or f"{model_name}_aug{int(augment)}"
    out = C.LOGS_DIR / f"cv_{name}.json"
    out.write_text(json.dumps(summary, indent=2))
    print(f"saved -> {out}")

    return summary
\end{lstlisting}

\subsection{\texttt{src/evaluate.py}}
Evaluation and embedding extraction.
\begin{lstlisting}[caption={Evaluation and embedding extraction (\texttt{src/evaluate.py})}]
"""Evaluation: aggregated cross-validation metrics and embedding extraction."""
from __future__ import annotations

import numpy as np
import torch
from torch.utils.data import DataLoader

from . import config as C
from .data import SpectrogramDataset, compute_norm_stats, make_fold_split
from .train import predict, train_one_fold


def evaluate_cv_predictions(X, meta, model_name="cnn", folds=None,
                            epochs=C.EPOCHS, augment=True, **kwargs):
    """Train per fold and pool the held-out predictions.

    Because every clip is in exactly one test fold, pooling gives one honest
    prediction for the entire dataset — the right basis for an aggregate
    confusion matrix and per-class scores.
    """
    import os
    import time

    folds = folds or list(range(1, C.N_FOLDS + 1))
    cooldown = float(os.environ.get("COOLDOWN_SEC", "0"))
    all_true, all_pred, all_prob, results = [], [], [], []

    for i, f in enumerate(folds):
        model, res = train_one_fold(
            X, meta, test_fold=f, model_name=model_name,
            epochs=epochs, augment=augment, **kwargs
        )
        results.append(res)

        (Xtr, ytr), _, (Xte, yte) = make_fold_split(X, meta, f)
        mean, std = compute_norm_stats(Xtr)
        dl = DataLoader(SpectrogramDataset(Xte, yte, mean, std), batch_size=C.BATCH_SIZE)
        yt, yp, pr = predict(model, dl, C.DEVICE)

        all_true.append(yt); all_pred.append(yp); all_prob.append(pr)

        if cooldown and i < len(folds) - 1:
            print(f"  [cooldown] pausing {cooldown:.0f}s (eval fold {i+1}/{len(folds)} done)")
            time.sleep(cooldown)

    return (np.concatenate(all_true), np.concatenate(all_pred),
            np.concatenate(all_prob), results)


def full_metrics(y_true, y_pred, class_names) -> dict:
    from sklearn.metrics import (accuracy_score, classification_report,
                                 confusion_matrix, f1_score,
                                 precision_score, recall_score)

    return {
        "accuracy": float(accuracy_score(y_true, y_pred)),
        "precision_macro": float(precision_score(y_true, y_pred, average="macro", zero_division=0)),
        "recall_macro": float(recall_score(y_true, y_pred, average="macro", zero_division=0)),
        "f1_macro": float(f1_score(y_true, y_pred, average="macro", zero_division=0)),
        "f1_weighted": float(f1_score(y_true, y_pred, average="weighted", zero_division=0)),
        "per_class_f1": {
            c: float(v) for c, v in zip(
                class_names, f1_score(y_true, y_pred, average=None, zero_division=0)
            )
        },
        "confusion_matrix": confusion_matrix(y_true, y_pred).tolist(),
        "report": classification_report(
            y_true, y_pred, target_names=class_names, zero_division=0, digits=3
        ),
    }


@torch.no_grad()
def extract_embeddings(model, X, y, mean, std, device=C.DEVICE, batch_size=C.BATCH_SIZE):
    """Pull the CNN's 256-d penultimate features for clustering / t-SNE."""
    model.eval()
    ds = SpectrogramDataset(X, y, mean, std, augment=False)
    dl = DataLoader(ds, batch_size=batch_size)

    embs, labels = [], []
    for xb, yb in dl:
        _, emb = model(xb.to(device), return_embedding=True)
        embs.append(emb.cpu().numpy())
        labels.append(yb.numpy())

    return np.concatenate(embs), np.concatenate(labels)
\end{lstlisting}

\subsection{\texttt{src/anomaly.py}}
Autoencoder novelty detection.
\begin{lstlisting}[caption={Autoencoder novelty detection (\texttt{src/anomaly.py})}]
"""Novel-sound detection with a convolutional autoencoder.

The classifier is closed-set: shown a smashing window it must still answer with
one of its ten labels, usually confidently and wrongly. For an assistive system
that is the dangerous failure mode.

The fix here is an autoencoder trained *only* on the ten known classes. It learns
to rebuild those, so an unfamiliar sound reconstructs badly and its reconstruction
error becomes a usable "I have not heard this before" score.
"""
from __future__ import annotations

import time

import numpy as np
import torch
import torch.nn as nn
from torch.utils.data import DataLoader

from . import config as C
from .data import SpectrogramDataset
from .models import ConvAutoencoder
from .train import set_seed


def train_autoencoder(X_train, mean, std, epochs=30, batch_size=C.BATCH_SIZE,
                      lr=1e-3, device=C.DEVICE, verbose=True, X_val=None):
    """Fit the autoencoder on known-class spectrograms only."""
    set_seed(C.SEED)

    dummy = np.zeros(len(X_train), dtype=np.int64)
    dl = DataLoader(SpectrogramDataset(X_train, dummy, mean, std, augment=False),
                    batch_size=batch_size, shuffle=True, drop_last=True)

    dl_val = None
    if X_val is not None:
        dl_val = DataLoader(
            SpectrogramDataset(X_val, np.zeros(len(X_val), dtype=np.int64), mean, std),
            batch_size=batch_size)

    model = ConvAutoencoder().to(device)
    criterion = nn.MSELoss()
    optimizer = torch.optim.AdamW(model.parameters(), lr=lr, weight_decay=1e-5)
    scheduler = torch.optim.lr_scheduler.CosineAnnealingLR(optimizer, T_max=epochs)

    history = {"train_loss": [], "val_loss": []}
    t0 = time.time()

    for ep in range(1, epochs + 1):
        model.train()
        total, n = 0.0, 0
        for xb, _ in dl:
            xb = xb.to(device)
            loss = criterion(model(xb), xb)
            optimizer.zero_grad(set_to_none=True)
            loss.backward()
            optimizer.step()
            total += loss.item() * xb.size(0); n += xb.size(0)
        scheduler.step()
        history["train_loss"].append(total / n)

        if dl_val is not None:
            model.eval()
            vt, vn = 0.0, 0
            with torch.no_grad():
                for xb, _ in dl_val:
                    xb = xb.to(device)
                    vt += criterion(model(xb), xb).item() * xb.size(0); vn += xb.size(0)
            history["val_loss"].append(vt / vn)

        if verbose and (ep % 5 == 0 or ep == 1):
            msg = f"  AE ep{ep:3d}/{epochs} train {history['train_loss'][-1]:.5f}"
            if history["val_loss"]:
                msg += f"  val {history['val_loss'][-1]:.5f}"
            print(msg)

    if verbose:
        print(f"  AE trained in {time.time()-t0:.0f}s")
    return model, history


@torch.no_grad()
def reconstruction_scores(model, X, mean, std, device=C.DEVICE,
                          batch_size=C.BATCH_SIZE) -> np.ndarray:
    """Per-clip reconstruction error — higher means less familiar."""
    model.eval()
    dl = DataLoader(SpectrogramDataset(X, np.zeros(len(X), dtype=np.int64), mean, std),
                    batch_size=batch_size)
    out = []
    for xb, _ in dl:
        out.append(model.reconstruction_error(xb.to(device)).cpu().numpy())
    return np.concatenate(out)


def evaluate_novelty(known_scores: np.ndarray, novel_scores: np.ndarray,
                     percentile: float = 95.0) -> dict:
    """Score the detector as a binary problem: novel (1) vs known (0).

    The threshold is set at a percentile of the KNOWN scores, so it can be
    calibrated without ever seeing a novel example — which is the situation a
    deployed system is actually in.
    """
    from sklearn.metrics import (average_precision_score, roc_auc_score, roc_curve)

    y = np.concatenate([np.zeros(len(known_scores)), np.ones(len(novel_scores))])
    s = np.concatenate([known_scores, novel_scores])

    auc = float(roc_auc_score(y, s))
    ap = float(average_precision_score(y, s))
    fpr, tpr, _ = roc_curve(y, s)

    thr = float(np.percentile(known_scores, percentile))
    tp = int((novel_scores > thr).sum())
    fn = int((novel_scores <= thr).sum())
    fp = int((known_scores > thr).sum())
    tn = int((known_scores <= thr).sum())

    precision = tp / (tp + fp) if (tp + fp) else 0.0
    recall = tp / (tp + fn) if (tp + fn) else 0.0
    f1 = 2 * precision * recall / (precision + recall) if (precision + recall) else 0.0

    return {
        "roc_auc": auc,
        "average_precision": ap,
        "threshold": thr,
        "threshold_percentile": percentile,
        "precision": precision,
        "recall": recall,
        "f1": f1,
        "confusion": {"tp": tp, "fp": fp, "tn": tn, "fn": fn},
        "fpr": fpr.tolist(),
        "tpr": tpr.tolist(),
        "known_mean": float(known_scores.mean()),
        "novel_mean": float(novel_scores.mean()),
    }


def softmax_confidence_baseline(probs_known: np.ndarray,
                                probs_novel: np.ndarray) -> dict:
    """Baseline: use 1 - max(softmax) as the novelty score.

    Included so the report can show whether the autoencoder actually beats the
    obvious "the classifier is unsure" heuristic, rather than assuming it does.
    """
    from sklearn.metrics import roc_auc_score

    known = 1.0 - probs_known.max(axis=1)
    novel = 1.0 - probs_novel.max(axis=1)
    y = np.concatenate([np.zeros(len(known)), np.ones(len(novel))])
    s = np.concatenate([known, novel])
    return {
        "roc_auc": float(roc_auc_score(y, s)),
        "known_mean_conf": float(probs_known.max(axis=1).mean()),
        "novel_mean_conf": float(probs_novel.max(axis=1).mean()),
    }
\end{lstlisting}

\subsection{\texttt{src/cluster.py}}
k-means and t-SNE clustering.
\begin{lstlisting}[caption={k-means and t-SNE clustering (\texttt{src/cluster.py})}]
"""Unsupervised analysis of the CNN's learned representation.

The classifier is trained with labels, but the *embedding space* it builds can be
examined without them. If k-means on those embeddings recovers the sound classes,
the network has learned a genuinely structured representation rather than a bare
decision rule — and the classes it confuses show up as overlapping clusters.
"""
from __future__ import annotations

import numpy as np

from . import config as C


def run_tsne(embeddings: np.ndarray, n_components: int = 2,
             perplexity: float = 30.0, max_samples: int = 3000,
             seed: int = C.SEED):
    """Project embeddings to 2-D for visualisation.

    t-SNE is O(n^2)-ish, so large inputs are subsampled. Returns (points, index)
    where `index` selects the rows that were kept.
    """
    from sklearn.manifold import TSNE

    rng = np.random.default_rng(seed)
    idx = (rng.choice(len(embeddings), max_samples, replace=False)
           if len(embeddings) > max_samples else np.arange(len(embeddings)))

    tsne = TSNE(n_components=n_components, perplexity=perplexity,
                init="pca", learning_rate="auto", random_state=seed)
    return tsne.fit_transform(embeddings[idx]), idx


def run_kmeans(embeddings: np.ndarray, n_clusters: int = 10, seed: int = C.SEED):
    from sklearn.cluster import KMeans

    km = KMeans(n_clusters=n_clusters, random_state=seed, n_init=10)
    return km.fit_predict(embeddings), km


def clustering_metrics(embeddings: np.ndarray, cluster_labels: np.ndarray,
                       true_labels: np.ndarray) -> dict:
    """Compare discovered clusters against the true classes.

    ARI/NMI are label-permutation invariant, which matters because cluster 3 has
    no reason to correspond to class 3. Silhouette judges cluster separation
    without using labels at all.
    """
    from sklearn.metrics import (adjusted_rand_score,
                                 normalized_mutual_info_score, silhouette_score)

    return {
        "adjusted_rand_index": float(adjusted_rand_score(true_labels, cluster_labels)),
        "normalized_mutual_info": float(normalized_mutual_info_score(true_labels, cluster_labels)),
        "silhouette": float(silhouette_score(embeddings, cluster_labels)),
        "n_clusters": int(len(np.unique(cluster_labels))),
    }


def cluster_purity(cluster_labels: np.ndarray, true_labels: np.ndarray,
                   class_names: list[str]) -> dict:
    """For each cluster, the dominant true class and how pure the cluster is."""
    out = {}
    for c in np.unique(cluster_labels):
        m = cluster_labels == c
        vals, counts = np.unique(true_labels[m], return_counts=True)
        top = vals[counts.argmax()]
        out[int(c)] = {
            "dominant_class": class_names[int(top)],
            "purity": float(counts.max() / counts.sum()),
            "size": int(m.sum()),
        }
    return out
\end{lstlisting}

\subsection{\texttt{src/viz.py}}
Figure generation.
\begin{lstlisting}[caption={Figure generation (\texttt{src/viz.py})}]
"""Figure generation for the report.

All figures share one visual system: a validated categorical palette, a single
blue sequential ramp for magnitude, recessive grid/axis ink, and direct labels
so identity never rests on colour alone.
"""
from __future__ import annotations

import matplotlib as mpl
import matplotlib.pyplot as plt
import numpy as np

from . import config as C

# --------------------------------------------------------------------------
# Design tokens (validated: all-pairs CVD dE 9.2, normal-vision 24.0, light)
# --------------------------------------------------------------------------
SERIES = ["#2a78d6", "#eb6834", "#1baf7a", "#eda100", "#e87ba4",
          "#008300", "#4a3aa7", "#e34948", "#6da7ec", "#c98500"]

SEQ_BLUE = ["#cde2fb", "#b7d3f6", "#9ec5f4", "#86b6ef", "#6da7ec",
            "#5598e7", "#3987e5", "#2a78d6", "#256abf", "#1c5cab",
            "#184f95", "#104281", "#0d366b"]

SURFACE = "#fcfcfb"
INK = "#0b0b0b"
INK_2 = "#52514e"
MUTED = "#898781"
GRID = "#e1e0d9"
AXIS = "#c3c2b7"

STATUS = {"good": "#0ca30c", "warning": "#fab219", "critical": "#d03b3b"}


def use_report_style() -> None:
    """Apply the shared style. Call once before generating figures."""
    mpl.rcParams.update({
        "figure.facecolor": SURFACE,
        "axes.facecolor": SURFACE,
        "savefig.facecolor": SURFACE,
        "font.family": "sans-serif",
        "font.sans-serif": ["Helvetica Neue", "Helvetica", "Arial", "DejaVu Sans"],
        "font.size": 10,
        "axes.titlesize": 12,
        "axes.titleweight": "bold",
        "axes.labelsize": 10,
        "axes.labelcolor": INK_2,
        "axes.edgecolor": AXIS,
        "axes.linewidth": 0.8,
        "axes.grid": True,
        "axes.axisbelow": True,
        "grid.color": GRID,
        "grid.linewidth": 0.7,
        "xtick.color": MUTED,
        "ytick.color": MUTED,
        "xtick.labelsize": 9,
        "ytick.labelsize": 9,
        "text.color": INK,
        "legend.frameon": False,
        "legend.fontsize": 9,
        "figure.dpi": 130,
        "savefig.dpi": 300,          # print-resolution figures for the report
        "savefig.bbox": "tight",
    })


def _despine(ax, keep=("left", "bottom")):
    for side in ("top", "right", "left", "bottom"):
        ax.spines[side].set_visible(side in keep)


def save(fig, name: str) -> str:
    """Save a figure to results/figures as PNG (report) and PDF (vector)."""
    path = C.FIGURES_DIR / f"{name}.png"
    fig.savefig(path)
    fig.savefig(C.FIGURES_DIR / f"{name}.pdf")
    plt.close(fig)
    print(f"  figure -> {path}")
    return str(path)


# --------------------------------------------------------------------------
# Dataset overview
# --------------------------------------------------------------------------
def plot_class_distribution(counts: dict[str, int], name="fig_class_distribution",
                            title="UrbanSound8K clips per class"):
    labels = list(counts.keys())
    values = [counts[k] for k in labels]
    order = np.argsort(values)[::-1]
    labels = [labels[i].replace("_", " ") for i in order]
    values = [values[i] for i in order]

    fig, ax = plt.subplots(figsize=(7.5, 4.2))
    bars = ax.barh(labels[::-1], values[::-1], color=SERIES[0], height=0.62)
    for b, v in zip(bars, values[::-1]):
        ax.text(v + max(values) * 0.012, b.get_y() + b.get_height() / 2,
                f"{v:,}", va="center", ha="left", fontsize=9, color=INK_2)

    ax.set_xlabel("Number of clips")
    ax.set_title(title, loc="left", pad=12)
    ax.set_xlim(0, max(values) * 1.12)
    ax.grid(axis="y", visible=False)
    _despine(ax)
    return save(fig, name)


def plot_spectrogram_grid(specs, titles, name="fig_spectrogram_examples",
                          suptitle="Log-mel spectrograms: each sound has a distinct signature",
                          ncols=5):
    n = len(specs)
    nrows = int(np.ceil(n / ncols))
    fig, axes = plt.subplots(nrows, ncols, figsize=(2.6 * ncols, 2.3 * nrows))
    axes = np.atleast_1d(axes).ravel()

    for ax, spec, t in zip(axes, specs, titles):
        ax.imshow(spec, origin="lower", aspect="auto", cmap="magma",
                  extent=[0, C.CLIP_SECONDS, 0, C.N_MELS])
        ax.set_title(t.replace("_", " "), fontsize=9.5, pad=5)
        ax.set_xticks([0, 2, 4]); ax.set_yticks([])
        ax.grid(False)
        ax.tick_params(labelsize=8)
    for ax in axes[n:]:
        ax.axis("off")

    fig.suptitle(suptitle, fontsize=12, fontweight="bold", x=0.02, ha="left", y=1.0)
    fig.supxlabel("Time (s)", fontsize=9, color=INK_2)
    fig.supylabel("Mel frequency band", fontsize=9, color=INK_2)
    fig.tight_layout()
    return save(fig, name)


# --------------------------------------------------------------------------
# Training
# --------------------------------------------------------------------------
def plot_training_curves(histories: dict[str, dict], name="fig_training_curves"):
    """histories: {model_label: history_dict}"""
    fig, (ax1, ax2) = plt.subplots(1, 2, figsize=(10, 3.8))

    for i, (label, h) in enumerate(histories.items()):
        c = SERIES[i]
        ep = range(1, len(h["train_loss"]) + 1)
        ax1.plot(ep, h["train_loss"], color=c, lw=2, ls="--", alpha=0.55)
        ax1.plot(ep, h["val_loss"], color=c, lw=2, label=label)
        ax2.plot(ep, [a * 100 for a in h["train_acc"]], color=c, lw=2, ls="--", alpha=0.55)
        ax2.plot(ep, [a * 100 for a in h["val_acc"]], color=c, lw=2, label=label)
        # direct label at the end of the validation curve
        ax2.annotate(label, (len(h["val_acc"]), h["val_acc"][-1] * 100),
                     xytext=(4, 0), textcoords="offset points",
                     color=c, fontsize=9, fontweight="bold", va="center")

    ax1.set_xlabel("Epoch"); ax1.set_ylabel("Cross-entropy loss")
    ax1.set_title("Loss  (dashed = train, solid = validation)", loc="left", pad=10)
    ax2.set_xlabel("Epoch"); ax2.set_ylabel("Accuracy (%)")
    ax2.set_title("Accuracy  (dashed = train, solid = validation)", loc="left", pad=10)
    ax1.legend(loc="upper right")
    for ax in (ax1, ax2):
        _despine(ax)
    fig.tight_layout()
    return save(fig, name)


def plot_confusion_matrix(cm, classes, name="fig_confusion_matrix",
                          title="CNN confusion matrix (10-fold aggregate)", normalize=True):
    from matplotlib.colors import LinearSegmentedColormap

    cmap = LinearSegmentedColormap.from_list("seqblue", SEQ_BLUE)
    m = cm.astype(float)
    if normalize:
        m = m / np.clip(m.sum(axis=1, keepdims=True), 1, None) * 100

    fig, ax = plt.subplots(figsize=(7.6, 6.4))
    im = ax.imshow(m, cmap=cmap, vmin=0, vmax=100 if normalize else m.max())

    labels = [c.replace("_", " ") for c in classes]
    ax.set_xticks(range(len(classes)), labels, rotation=45, ha="right")
    ax.set_yticks(range(len(classes)), labels)
    ax.set_xlabel("Predicted"); ax.set_ylabel("Actual")
    ax.set_title(title, loc="left", pad=12)
    ax.grid(False)

    # value labels — identity never rests on colour alone
    for i in range(m.shape[0]):
        for j in range(m.shape[1]):
            v = m[i, j]
            if v >= 0.5:
                ax.text(j, i, f"{v:.0f}", ha="center", va="center", fontsize=8.5,
                        color="#ffffff" if v > 55 else INK_2,
                        fontweight="bold" if i == j else "normal")

    cb = fig.colorbar(im, ax=ax, fraction=0.045, pad=0.03)
    cb.set_label("% of actual class" if normalize else "clips", color=INK_2, fontsize=9)
    cb.outline.set_visible(False)
    fig.tight_layout()
    return save(fig, name)


def plot_model_comparison(results: dict[str, dict], name="fig_model_comparison",
                          n_folds: int | None = None):
    """results: {label: {'acc_mean','acc_std','f1_mean','f1_std','n_params'}}"""
    labels = list(results.keys())
    x = np.arange(len(labels))
    w = 0.36

    fig, (ax1, ax2) = plt.subplots(1, 2, figsize=(10.5, 4.3),
                                   gridspec_kw={"width_ratios": [1.55, 1]})

    acc = [results[l]["acc_mean"] * 100 for l in labels]
    accs = [results[l]["acc_std"] * 100 for l in labels]
    f1 = [results[l]["f1_mean"] * 100 for l in labels]
    f1s = [results[l]["f1_std"] * 100 for l in labels]

    b1 = ax1.bar(x - w / 2, acc, w, yerr=accs, capsize=3, color=SERIES[0], label="Accuracy")
    b2 = ax1.bar(x + w / 2, f1, w, yerr=f1s, capsize=3, color=SERIES[1], label="Macro F1")

    headroom = max(v + e for v, e in zip(acc + f1, accs + f1s))
    # Labels clear the error-bar cap, not just the bar top.
    for bars, vals, errs in ((b1, acc, accs), (b2, f1, f1s)):
        for b, v, e in zip(bars, vals, errs):
            ax1.text(b.get_x() + b.get_width() / 2, v + e + headroom * 0.035, f"{v:.1f}",
                     ha="center", fontsize=9, color=INK_2, fontweight="bold")

    fold_note = f"{n_folds}-fold " if n_folds else ""
    ax1.set_xticks(x, labels)
    ax1.set_ylabel("Score (%)")
    ax1.set_ylim(0, headroom * 1.30)
    ax1.set_title(f"Performance ({fold_note}mean ± sd)", loc="left", pad=10)
    # Legend on the title line but right-aligned: clear of both the
    # left-aligned title and every bar.
    ax1.legend(loc="lower right", bbox_to_anchor=(1.0, 1.005), ncols=2,
               frameon=False, borderaxespad=0)
    ax1.grid(axis="x", visible=False)
    _despine(ax1)

    params = [results[l]["n_params"] / 1e6 for l in labels]
    b3 = ax2.bar(x, params, 0.5, color=SERIES[2])
    for b, v in zip(b3, params):
        ax2.text(b.get_x() + b.get_width() / 2, v + max(params) * 0.03,
                 f"{v:.2f}M", ha="center", fontsize=9, color=INK_2, fontweight="bold")
    ax2.set_xticks(x, labels)
    ax2.set_ylabel("Trainable parameters (millions)")
    ax2.set_ylim(0, max(params) * 1.22)
    ax2.set_title("Model size", loc="left", pad=10)
    ax2.grid(axis="x", visible=False)
    _despine(ax2)

    fig.tight_layout()
    return save(fig, name)


def plot_per_class_f1(per_class: dict[str, float], name="fig_per_class_f1"):
    items = sorted(per_class.items(), key=lambda kv: kv[1])
    labels = [k.replace("_", " ") for k, _ in items]
    vals = [v * 100 for _, v in items]
    colors = [STATUS["critical"] if v < 70 else (STATUS["warning"] if v < 85 else SERIES[0])
              for v in vals]

    fig, ax = plt.subplots(figsize=(7.5, 4.2))
    bars = ax.barh(labels, vals, color=colors, height=0.62)
    for b, v in zip(bars, vals):
        ax.text(v + 1, b.get_y() + b.get_height() / 2, f"{v:.1f}",
                va="center", fontsize=9, color=INK_2)
    ax.set_xlabel("F1 score (%)")
    ax.set_xlim(0, 105)
    ax.set_title("Per-class F1 — which sounds the CNN finds hard", loc="left", pad=12)
    ax.grid(axis="y", visible=False)
    _despine(ax)
    return save(fig, name)


# --------------------------------------------------------------------------
# Embeddings / clustering
# --------------------------------------------------------------------------
def plot_tsne(emb2d, labels, class_names, name="fig_tsne_embeddings",
              title="What the CNN learned: t-SNE of 256-d embeddings"):
    """Colour + a direct centroid label per class, so identity is never
    carried by colour alone (10 classes exceeds a safe colour-only cap)."""
    fig, ax = plt.subplots(figsize=(8.2, 6.8))

    for i, cname in enumerate(class_names):
        m = labels == i
        if not m.any():
            continue
        ax.scatter(emb2d[m, 0], emb2d[m, 1], s=9, alpha=0.55,
                   color=SERIES[i % len(SERIES)], linewidths=0)
        cx, cy = emb2d[m, 0].mean(), emb2d[m, 1].mean()
        ax.annotate(cname.replace("_", " "), (cx, cy), fontsize=9.5, fontweight="bold",
                    ha="center", va="center", color=INK,
                    bbox=dict(boxstyle="round,pad=0.28", fc=SURFACE,
                              ec=SERIES[i % len(SERIES)], lw=1.4, alpha=0.92))

    ax.set_xticks([]); ax.set_yticks([])
    ax.set_title(title, loc="left", pad=12)
    ax.grid(False)
    _despine(ax, keep=())
    fig.tight_layout()
    return save(fig, name)


# --------------------------------------------------------------------------
# Anomaly / novelty detection
# --------------------------------------------------------------------------
def plot_anomaly_scores(known, novel, threshold=None, name="fig_anomaly_scores"):
    fig, ax = plt.subplots(figsize=(8, 4.2))
    bins = np.linspace(0, np.percentile(np.concatenate([known, novel]), 99), 60)

    ax.hist(known, bins=bins, color=SERIES[0], alpha=0.75, label="Known sounds (UrbanSound8K)")
    ax.hist(novel, bins=bins, color=SERIES[1], alpha=0.75, label="Novel sounds (ESC-50, unseen)")

    if threshold is not None:
        ax.axvline(threshold, color=STATUS["critical"], lw=2, ls="--")
        ax.annotate(f"threshold {threshold:.3f}", (threshold, ax.get_ylim()[1] * 0.92),
                    xytext=(7, 0), textcoords="offset points",
                    color=STATUS["critical"], fontsize=9, fontweight="bold")

    ax.set_xlabel("Autoencoder reconstruction error")
    ax.set_ylabel("Number of clips")
    ax.set_title("Novelty detection: unseen sounds reconstruct worse", loc="left", pad=12)
    ax.legend(loc="upper right")
    ax.grid(axis="x", visible=False)
    _despine(ax)
    return save(fig, name)


def plot_roc(fpr, tpr, auc, name="fig_anomaly_roc"):
    fig, ax = plt.subplots(figsize=(5.2, 4.8))
    ax.plot([0, 1], [0, 1], color=AXIS, lw=1.4, ls="--")
    ax.plot(fpr, tpr, color=SERIES[0], lw=2.4)
    ax.annotate(f"AUC = {auc:.3f}", (0.55, 0.22), fontsize=12, fontweight="bold", color=SERIES[0])
    ax.set_xlabel("False positive rate"); ax.set_ylabel("True positive rate")
    ax.set_title("Novel-sound detection ROC", loc="left", pad=12)
    ax.set_xlim(0, 1); ax.set_ylim(0, 1.02)
    _despine(ax)
    fig.tight_layout()
    return save(fig, name)


def plot_reconstructions(originals, recons, titles, name="fig_ae_reconstructions"):
    n = len(originals)
    fig, axes = plt.subplots(2, n, figsize=(2.5 * n, 4.6))
    axes = np.atleast_2d(axes)
    for j in range(n):
        for row, (data, tag) in enumerate(((originals[j], "input"), (recons[j], "rebuilt"))):
            ax = axes[row, j]
            ax.imshow(data, origin="lower", aspect="auto", cmap="magma")
            ax.set_xticks([]); ax.set_yticks([]); ax.grid(False)
            if row == 0:
                ax.set_title(titles[j], fontsize=9.5, pad=5)
            if j == 0:
                ax.set_ylabel(tag, fontsize=9.5, color=INK_2)
    fig.suptitle("Autoencoder reconstructions — known sounds rebuild cleanly, novel ones do not",
                 fontsize=12, fontweight="bold", x=0.02, ha="left")
    fig.tight_layout()
    return save(fig, name)
\end{lstlisting}

\subsection{\texttt{run\_experiments.py}}
Staged experiment runner.
\begin{lstlisting}[caption={Staged experiment runner (\texttt{run\_experiments.py})}]
#!/usr/bin/env python
"""End-to-end experiment runner.

Stages (select with --stage, default: all)

  prep     extract and cache mel-spectrograms for both datasets
  explore  dataset figures (class balance, example spectrograms)
  compare  MLP vs CNN, and CNN with/without augmentation
  final    full 10-fold CNN run -> confusion matrix, per-class F1, saved model
  cluster  k-means + t-SNE over the CNN's learned embeddings
  anomaly  autoencoder novelty detection vs unseen ESC-50 sounds

Examples
  python run_experiments.py --stage prep
  python run_experiments.py --stage compare --folds 1 2 3 --epochs 25
  python run_experiments.py                      # everything, full protocol
"""
from __future__ import annotations

import argparse
import json
import time

import numpy as np
import torch

from src import config as C
from src import viz
from src.data import (load_esc50_metadata, load_us8k_metadata,
                      compute_norm_stats, make_fold_split, precompute_features)


def _save(name: str, obj) -> None:
    p = C.LOGS_DIR / f"{name}.json"
    p.write_text(json.dumps(obj, indent=2, default=str))
    print(f"  saved -> {p}")


def banner(msg: str) -> None:
    print(f"\n{'='*74}\n  {msg}\n{'='*74}")


# ---------------------------------------------------------------------------
def stage_prep(force=False):
    banner("STAGE: prep — extracting mel-spectrograms")

    us8k = load_us8k_metadata()
    print(f"UrbanSound8K: {len(us8k)} clips, {us8k['class'].nunique()} classes")
    X, meta = precompute_features(us8k, "us8k", force=force)

    esc = load_esc50_metadata()
    novel = esc[esc["category"].isin(C.ESC50_NOVEL_CLASSES)].reset_index(drop=True)
    print(f"ESC-50 novel subset: {len(novel)} clips, {novel['category'].nunique()} classes")
    Xn, mn = precompute_features(novel, "esc50_novel", force=force)

    _save("prep_summary", {
        "us8k_shape": list(X.shape),
        "us8k_class_counts": us8k["class"].value_counts().to_dict(),
        "esc50_novel_shape": list(Xn.shape),
        "esc50_novel_classes": sorted(novel["category"].unique().tolist()),
        "feature": {"sample_rate": C.SAMPLE_RATE, "n_mels": C.N_MELS,
                    "n_fft": C.N_FFT, "hop_length": C.HOP_LENGTH,
                    "clip_seconds": C.CLIP_SECONDS, "n_frames": C.N_FRAMES},
    })
    return X, meta


# ---------------------------------------------------------------------------
def stage_explore(X, meta):
    banner("STAGE: explore — dataset figures")
    viz.use_report_style()

    us8k = load_us8k_metadata()
    viz.plot_class_distribution(us8k["class"].value_counts().to_dict())

    # one representative spectrogram per class
    specs, titles = [], []
    for cid in range(10):
        rows = np.where(meta["classID"].values == cid)[0]
        if len(rows):
            specs.append(X[rows[0]])
            titles.append(C.US8K_CLASSES[cid])
    viz.plot_spectrogram_grid(specs, titles, ncols=5)

    # the same class under augmentation, to illustrate what the model sees
    from src.data import spec_augment, time_shift
    rng = np.random.default_rng(C.SEED)
    base = X[np.where(meta["classID"].values == 8)[0][0]]  # siren
    variants = [base, time_shift(base, rng), spec_augment(base, rng),
                spec_augment(time_shift(base, rng), rng)]
    viz.plot_spectrogram_grid(
        variants, ["original", "time shift", "SpecAugment", "both"],
        name="fig_augmentation", ncols=4,
        suptitle="Augmentation: the same siren, varied so the CNN cannot memorise it")


# ---------------------------------------------------------------------------
def stage_compare(X, meta, folds, epochs):
    banner("STAGE: compare — MLP vs CNN, and augmentation ablation")
    from src.train import cross_validate, train_one_fold

    results = {}
    results["MLP"] = cross_validate(X, meta, "mlp", folds=folds, epochs=epochs,
                                    augment=False, tag="mlp")
    results["CNN"] = cross_validate(X, meta, "cnn", folds=folds, epochs=epochs,
                                    augment=False, tag="cnn_noaug")
    results["CNN + augment"] = cross_validate(X, meta, "cnn", folds=folds, epochs=epochs,
                                              augment=True, tag="cnn_aug")

    viz.use_report_style()
    viz.plot_model_comparison(results, n_folds=len(folds))

    # training curves on a single representative fold
    f = folds[0]
    hists = {}
    for label, (mname, aug) in {
        "MLP": ("mlp", False),
        "CNN": ("cnn", False),
        "CNN + augment": ("cnn", True),
    }.items():
        _, res = train_one_fold(X, meta, test_fold=f, model_name=mname,
                                epochs=epochs, augment=aug, verbose=False)
        hists[label] = res.history
    viz.plot_training_curves(hists)

    _save("comparison_summary", results)
    return results


# ---------------------------------------------------------------------------
def stage_final(X, meta, folds, epochs):
    banner("STAGE: final — full CNN evaluation")
    from src.evaluate import evaluate_cv_predictions, full_metrics

    y_true, y_pred, y_prob, per_fold = evaluate_cv_predictions(
        X, meta, "cnn", folds=folds, epochs=epochs, augment=True)

    m = full_metrics(y_true, y_pred, C.US8K_CLASSES)
    print("\n" + m["report"])
    print(f"Accuracy {m['accuracy']:.4f} | macro-F1 {m['f1_macro']:.4f}")

    viz.use_report_style()
    viz.plot_confusion_matrix(np.array(m["confusion_matrix"]), C.US8K_CLASSES)
    viz.plot_per_class_f1(m["per_class_f1"])

    _save("final_metrics", {k: v for k, v in m.items() if k != "report"})
    (C.LOGS_DIR / "final_classification_report.txt").write_text(m["report"])

    # Train and persist one deployable model (held-out fold 10) for the demo.
    from src.train import train_one_fold
    banner("training deployable model for the live demo (test fold 10)")
    model, res = train_one_fold(X, meta, test_fold=10, model_name="cnn",
                                epochs=epochs, augment=True,
                                save_path=C.MODELS_DIR / "cnn_demo.pt")
    print(f"demo model saved: test acc {res.test_acc:.4f}")
    return m, model


# ---------------------------------------------------------------------------
def stage_cluster(X, meta, epochs):
    banner("STAGE: cluster — k-means + t-SNE on CNN embeddings")
    from src.cluster import (cluster_purity, clustering_metrics, run_kmeans, run_tsne)
    from src.evaluate import extract_embeddings
    from src.train import train_one_fold

    test_fold = 10
    model, _ = train_one_fold(X, meta, test_fold=test_fold, model_name="cnn",
                              epochs=epochs, augment=True, verbose=False)

    (Xtr, ytr), _, (Xte, yte) = make_fold_split(X, meta, test_fold)
    mean, std = compute_norm_stats(Xtr)

    # Embeddings come from the held-out fold — never data the CNN trained on.
    emb, labels = extract_embeddings(model, Xte, yte, mean, std)
    print(f"embeddings: {emb.shape}")

    clusters, _ = run_kmeans(emb, n_clusters=10)
    metrics = clustering_metrics(emb, clusters, labels)
    purity = cluster_purity(clusters, labels, C.US8K_CLASSES)
    print(f"ARI {metrics['adjusted_rand_index']:.3f} | "
          f"NMI {metrics['normalized_mutual_info']:.3f} | "
          f"silhouette {metrics['silhouette']:.3f}")

    pts, idx = run_tsne(emb)
    viz.use_report_style()
    viz.plot_tsne(pts, labels[idx], C.US8K_CLASSES)

    _save("cluster_metrics", {"metrics": metrics, "cluster_purity": purity})
    return metrics


# ---------------------------------------------------------------------------
def stage_anomaly(X, meta, epochs):
    banner("STAGE: anomaly — autoencoder novelty detection")
    from src.anomaly import (evaluate_novelty, reconstruction_scores,
                             softmax_confidence_baseline, train_autoencoder)
    from src.data import SpectrogramDataset
    from src.train import predict, train_one_fold
    from torch.utils.data import DataLoader

    Xn = np.load(C.PROCESSED_DIR / "esc50_novel_X.npy")
    print(f"novel (unseen) clips: {Xn.shape}")

    test_fold = 10
    (Xtr, ytr), (Xva, yva), (Xte, yte) = make_fold_split(X, meta, test_fold)
    mean, std = compute_norm_stats(Xtr)

    ae, hist = train_autoencoder(Xtr, mean, std, epochs=epochs, X_val=Xva)

    known = reconstruction_scores(ae, Xte, mean, std)
    novel = reconstruction_scores(ae, Xn, mean, std)
    res = evaluate_novelty(known, novel)
    print(f"AE novelty  ROC-AUC {res['roc_auc']:.3f} | AP {res['average_precision']:.3f} | "
          f"P {res['precision']:.3f} R {res['recall']:.3f} F1 {res['f1']:.3f}")

    # Baseline: does classifier under-confidence already do this job?
    clf, _ = train_one_fold(X, meta, test_fold=test_fold, model_name="cnn",
                            epochs=epochs, augment=True, verbose=False)
    dl_k = DataLoader(SpectrogramDataset(Xte, yte, mean, std), batch_size=C.BATCH_SIZE)
    dl_n = DataLoader(SpectrogramDataset(Xn, np.zeros(len(Xn), dtype=np.int64), mean, std),
                      batch_size=C.BATCH_SIZE)
    _, _, pk = predict(clf, dl_k, C.DEVICE)
    _, _, pn = predict(clf, dl_n, C.DEVICE)
    base = softmax_confidence_baseline(pk, pn)
    print(f"softmax baseline ROC-AUC {base['roc_auc']:.3f} "
          f"(mean confidence: known {base['known_mean_conf']:.3f}, "
          f"novel {base['novel_mean_conf']:.3f})")

    viz.use_report_style()
    viz.plot_anomaly_scores(known, novel, threshold=res["threshold"])
    viz.plot_roc(np.array(res["fpr"]), np.array(res["tpr"]), res["roc_auc"])

    # visual side-by-side of reconstructions
    with torch.no_grad():
        def recon(arr, i):
            x = torch.from_numpy((arr[i] - mean) / std).float()[None, None].to(C.DEVICE)
            return x[0, 0].cpu().numpy(), ae(x)[0, 0].cpu().numpy()

        pairs = [recon(Xte, 0), recon(Xte, 1), recon(Xn, 0), recon(Xn, 1)]
    viz.plot_reconstructions([p[0] for p in pairs], [p[1] for p in pairs],
                             ["known", "known", "novel", "novel"])

    torch.save({"state_dict": ae.state_dict(), "mean": mean, "std": std,
                "threshold": res["threshold"]}, C.MODELS_DIR / "autoencoder.pt")

    _save("anomaly_results", {
        "autoencoder": {k: v for k, v in res.items() if k not in ("fpr", "tpr")},
        "softmax_baseline": base,
        "ae_history": hist,
    })
    return res


# ---------------------------------------------------------------------------
def main():
    ap = argparse.ArgumentParser(description=__doc__,
                                 formatter_class=argparse.RawDescriptionHelpFormatter)
    ap.add_argument("--stage", default="all",
                    choices=["all", "prep", "explore", "compare", "final",
                             "cluster", "anomaly"])
    ap.add_argument("--folds", type=int, nargs="+", default=None,
                    help="folds to evaluate (default: all 10)")
    ap.add_argument("--epochs", type=int, default=C.EPOCHS)
    ap.add_argument("--quick", action="store_true",
                    help="fast smoke run: 1 fold, few epochs")
    ap.add_argument("--force-prep", action="store_true")
    args = ap.parse_args()

    folds = args.folds or list(range(1, C.N_FOLDS + 1))
    epochs = args.epochs
    if args.quick:
        folds, epochs = [1], 3
        print(">>> QUICK MODE: 1 fold, 3 epochs")

    print(f"device: {C.DEVICE} | torch {torch.__version__}")
    print(f"folds: {folds} | epochs: {epochs}")
    t0 = time.time()

    X, meta = stage_prep(force=args.force_prep)
    s = args.stage

    if s in ("all", "explore"):
        stage_explore(X, meta)
    if s in ("all", "compare"):
        stage_compare(X, meta, folds, epochs)
    if s in ("all", "final"):
        stage_final(X, meta, folds, epochs)
    if s in ("all", "cluster"):
        stage_cluster(X, meta, epochs)
    if s in ("all", "anomaly"):
        stage_anomaly(X, meta, epochs)

    banner(f"DONE in {(time.time()-t0)/60:.1f} min — figures in {C.FIGURES_DIR}")


if __name__ == "__main__":
    main()
\end{lstlisting}

\subsection{\texttt{demo/live\_demo.py}}
Real-time microphone demo.
\begin{lstlisting}[caption={Real-time microphone demo (\texttt{demo/live\_demo.py})}]
#!/usr/bin/env python
"""Real-time acoustic scene awareness from the laptop microphone.

Holds a rolling 4-second window of audio, converts it to a log-mel spectrogram
on the same front-end used in training, and runs both networks each refresh:

  * the CNN classifies the sound and drives the confidence bars
  * the autoencoder scores how familiar it is, so an unrecognised sound is
    reported as UNKNOWN rather than forced into one of the ten labels

Run:
    python demo/live_demo.py                 # live microphone
    python demo/live_demo.py --list-devices
    python demo/live_demo.py --file clip.wav # replay a file instead
"""
from __future__ import annotations

import argparse
import queue
import sys
import threading
from pathlib import Path

import matplotlib
import matplotlib.pyplot as plt
import numpy as np
import torch

sys.path.insert(0, str(Path(__file__).resolve().parent.parent))

from src import config as C  # noqa: E402
from src import viz  # noqa: E402
from src.data import audio_to_melspec  # noqa: E402
from src.models import AudioCNN, ConvAutoencoder  # noqa: E402

REFRESH_SEC = 0.5          # how often inference runs
SMOOTHING = 0.6            # EMA over probabilities, damps single-frame flicker


# ---------------------------------------------------------------------------
class Engine:
    """Loads the trained models and turns a waveform into a prediction."""

    def __init__(self, cnn_path, ae_path=None, device=C.DEVICE):
        self.device = device

        ck = torch.load(cnn_path, map_location=device, weights_only=False)
        self.mean, self.std = ck["mean"], ck["std"]
        self.cnn = AudioCNN(n_classes=ck.get("n_classes", 10)).to(device)
        self.cnn.load_state_dict(ck["state_dict"])
        self.cnn.eval()

        self.ae = None
        self.threshold = None
        if ae_path and Path(ae_path).exists():
            ack = torch.load(ae_path, map_location=device, weights_only=False)
            self.ae = ConvAutoencoder().to(device)
            self.ae.load_state_dict(ack["state_dict"])
            self.ae.eval()
            self.threshold = ack.get("threshold")

        self.probs = np.zeros(len(C.US8K_CLASSES))

    @torch.no_grad()
    def infer(self, wave: np.ndarray):
        spec = audio_to_melspec(wave)
        x = torch.from_numpy((spec - self.mean) / self.std).float()[None, None].to(self.device)

        p = torch.softmax(self.cnn(x), dim=1)[0].cpu().numpy()
        self.probs = SMOOTHING * self.probs + (1 - SMOOTHING) * p  # EMA

        novelty, is_novel = None, False
        if self.ae is not None:
            novelty = float(self.ae.reconstruction_error(x).item())
            if self.threshold is not None:
                is_novel = novelty > self.threshold

        return spec, self.probs, novelty, is_novel


# ---------------------------------------------------------------------------
class AudioStream:
    """Rolling microphone buffer of the last CLIP_SECONDS seconds."""

    def __init__(self, device=None, samplerate=C.SAMPLE_RATE):
        self.samplerate = samplerate
        self.buffer = np.zeros(C.N_SAMPLES, dtype=np.float32)
        self.lock = threading.Lock()
        self.q: queue.Queue = queue.Queue()
        self.device = device
        self.stream = None
        self.level = 0.0

    def _callback(self, indata, frames, time_info, status):
        if status:
            print(f"[audio] {status}", file=sys.stderr)
        mono = indata[:, 0].copy()
        with self.lock:
            n = len(mono)
            self.buffer = np.roll(self.buffer, -n)
            self.buffer[-n:] = mono
            self.level = float(np.abs(mono).max())

    def start(self):
        import sounddevice as sd

        self.stream = sd.InputStream(
            channels=1, samplerate=self.samplerate, device=self.device,
            blocksize=int(self.samplerate * 0.1), callback=self._callback,
        )
        self.stream.start()
        print(f"[audio] listening at {self.samplerate} Hz")

    def read(self) -> np.ndarray:
        with self.lock:
            return self.buffer.copy()

    def stop(self):
        if self.stream:
            self.stream.stop(); self.stream.close()


# ---------------------------------------------------------------------------
class Dashboard:
    """Live figure: spectrogram, confidence bars, and a status banner."""

    def __init__(self, engine: Engine, source, is_file=False):
        self.engine = engine
        self.source = source
        self.is_file = is_file
        viz.use_report_style()

        self.fig = plt.figure(figsize=(13, 6.6))
        self.fig.canvas.manager.set_window_title("Acoustic Scene Awareness — live")
        gs = self.fig.add_gridspec(2, 2, height_ratios=[1, 5.2],
                                   width_ratios=[1.35, 1], hspace=0.28, wspace=0.22)

        # Status banner
        self.ax_banner = self.fig.add_subplot(gs[0, :])
        self.ax_banner.axis("off")
        self.banner = self.ax_banner.text(
            0.5, 0.5, "listening…", ha="center", va="center",
            fontsize=25, fontweight="bold", color=viz.INK,
            transform=self.ax_banner.transAxes,
            bbox=dict(boxstyle="round,pad=0.55", fc="#eef3fa", ec=viz.SERIES[0], lw=2.5))

        # Spectrogram
        self.ax_spec = self.fig.add_subplot(gs[1, 0])
        self.im = self.ax_spec.imshow(
            np.full((C.N_MELS, C.N_FRAMES), -80.0), origin="lower", aspect="auto",
            cmap="magma", vmin=-80, vmax=0, extent=[0, C.CLIP_SECONDS, 0, C.N_MELS])
        self.ax_spec.set_title("Live log-mel spectrogram (last 4 s)", loc="left", pad=10)
        self.ax_spec.set_xlabel("Time (s)"); self.ax_spec.set_ylabel("Mel band")
        self.ax_spec.grid(False)

        # Confidence bars
        self.ax_bars = self.fig.add_subplot(gs[1, 1])
        labels = [C.DEMO_LABELS[c] for c in C.US8K_CLASSES]
        self.bars = self.ax_bars.barh(labels, np.zeros(len(labels)),
                                      color=viz.SERIES[0], height=0.62)
        self.val_txt = [
            self.ax_bars.text(0.6, b.get_y() + b.get_height() / 2, "", va="center",
                              fontsize=9, color=viz.INK_2)
            for b in self.bars
        ]
        self.ax_bars.set_xlim(0, 1)
        self.ax_bars.set_xlabel("Confidence")
        self.ax_bars.set_title("CNN class confidence", loc="left", pad=10)
        self.ax_bars.grid(axis="y", visible=False)
        for s in ("top", "right"):
            self.ax_bars.spines[s].set_visible(False)

        self.file_pos = 0

    def _next_wave(self) -> np.ndarray:
        if not self.is_file:
            return self.source.read()
        # File mode: sweep a window across the clip so the demo animates.
        wave = self.source
        step = int(C.SAMPLE_RATE * REFRESH_SEC)
        start = self.file_pos % max(len(wave), 1)
        self.file_pos += step
        w = wave[start:start + C.N_SAMPLES]
        if len(w) < C.N_SAMPLES:
            w = np.pad(w, (0, C.N_SAMPLES - len(w)))
        return w

    def update(self, _frame):
        wave = self._next_wave()

        if np.abs(wave).max() < 1e-4:                     # effectively silence
            self.banner.set_text("listening…")
            self.banner.set_color(viz.INK)
            self.banner.get_bbox_patch().set(fc="#f2f2ef", ec=viz.AXIS)
            for b, t in zip(self.bars, self.val_txt):
                b.set_width(0); t.set_text("")
            return

        spec, probs, novelty, is_novel = self.engine.infer(wave)
        self.im.set_data(spec)

        top = int(np.argmax(probs))
        conf = float(probs[top])
        cls = C.US8K_CLASSES[top]

        for b, t, p in zip(self.bars, self.val_txt, probs):
            b.set_width(p)
            b.set_color(viz.SERIES[0] if p < 0.5 else viz.SERIES[1])
            t.set_x(min(p + 0.02, 0.88))
            t.set_text(f"{p*100:.0f}%" if p > 0.02 else "")

        if is_novel:
            txt, fc, ec, col = ("❓ UNKNOWN SOUND", "#fdf3e3",
                                viz.STATUS["warning"], "#7a5200")
        elif cls in C.ALERT_CLASSES and conf > 0.45:
            txt = f"⚠️  {C.DEMO_LABELS[cls].upper()}"
            fc, ec, col = "#fdeaea", viz.STATUS["critical"], "#8c1d1d"
        elif conf > 0.35:
            txt, fc, ec, col = (C.DEMO_LABELS[cls], "#eef3fa", viz.SERIES[0], viz.INK)
        else:
            txt, fc, ec, col = ("listening…", "#f2f2ef", viz.AXIS, viz.INK)

        if conf > 0.35 and not is_novel:
            txt += f"   {conf*100:.0f}%"
        if novelty is not None:
            txt += f"      novelty {novelty:.3f}"

        self.banner.set_text(txt)
        self.banner.set_color(col)
        self.banner.get_bbox_patch().set(fc=fc, ec=ec)

    def run(self):
        from matplotlib.animation import FuncAnimation

        self.anim = FuncAnimation(self.fig, self.update,
                                  interval=int(REFRESH_SEC * 1000),
                                  cache_frame_data=False)
        plt.tight_layout()
        plt.show()


# ---------------------------------------------------------------------------
def main():
    ap = argparse.ArgumentParser(description=__doc__,
                                 formatter_class=argparse.RawDescriptionHelpFormatter)
    ap.add_argument("--model", default=str(C.MODELS_DIR / "cnn_demo.pt"))
    ap.add_argument("--autoencoder", default=str(C.MODELS_DIR / "autoencoder.pt"))
    ap.add_argument("--device", type=int, default=None, help="input device index")
    ap.add_argument("--file", default=None, help="replay a wav file instead of the mic")
    ap.add_argument("--list-devices", action="store_true")
    args = ap.parse_args()

    if args.list_devices:
        import sounddevice as sd
        print(sd.query_devices())
        return

    if not Path(args.model).exists():
        sys.exit(f"Model not found: {args.model}\n"
                 f"Train one first:  python run_experiments.py --stage final")

    # pyplot is already imported (via src.viz), so the backend has to be
    # switched rather than set — matplotlib.use() would be a no-op here.
    for backend in (("MacOSX", "TkAgg") if sys.platform == "darwin" else ("TkAgg",)):
        try:
            plt.switch_backend(backend)
            break
        except Exception:
            continue
    if matplotlib.get_backend().lower() in ("agg", "template"):
        sys.exit("No interactive matplotlib backend available — the live demo "
                 "needs a display. Try: uv pip install pyqt6")
    print(f"[ui] backend: {matplotlib.get_backend()}")

    engine = Engine(args.model, args.autoencoder)
    print(f"[model] loaded on {engine.device}"
          f"{' (+ autoencoder)' if engine.ae is not None else ''}")

    if args.file:
        import librosa
        wave, _ = librosa.load(args.file, sr=C.SAMPLE_RATE, mono=True)
        print(f"[file] {args.file} — {len(wave)/C.SAMPLE_RATE:.1f}s")
        Dashboard(engine, wave, is_file=True).run()
        return

    stream = AudioStream(device=args.device)
    stream.start()
    try:
        Dashboard(engine, stream).run()
    finally:
        stream.stop()
        print("[audio] stopped")


if __name__ == "__main__":
    main()
\end{lstlisting}
